\section{Background}
Railway interlocking is a critical safety mechanism that ensures the secure movement of trains through junctions, crossings, and stations by preventing conflicting routes from being set simultaneously. Historically, this function has been carried out through mechanical frames and relay-based logic systems—robust but inflexible solutions that are difficult to scale or adapt. As rail infrastructure becomes more complex and digital transformation accelerates across industries, these traditional systems are increasingly becoming a bottleneck to achieving high-performance, safe, and automated train operations.

With the rising demand for intelligent transport systems, the railway industry is witnessing a paradigm shift toward software-based interlocking, where logic is defined through configurable rule engines and verified in real time. These systems offer greater flexibility, ease of maintenance, and dynamic adaptability, making them ideal for modern urban networks and developing regions upgrading from legacy infrastructure.
